\documentclass[12pt,a4paper]{report}

% Packages
\usepackage[utf8]{inputenc}
\usepackage[T1]{fontenc}
\usepackage[lining]{ebgaramond}
\usepackage[cmintegrals,cmbraces]{newtxmath}
\usepackage{microtype}
\usepackage{geometry}
\usepackage{tabularx}
\usepackage{graphicx}
\usepackage{longtable}
\usepackage{booktabs}
\usepackage{array}
\usepackage{listings}
\usepackage{xcolor}
\usepackage{hyperref}
\usepackage{amsmath}
\usepackage{pdflscape}
\usepackage{float}
\usepackage{enumitem}
\graphicspath{{./}}

% Page layout
\geometry{
    left=1.8cm,
    right=1.8cm,
    top=2.2cm,
    bottom=2.2cm
}

% Color definitions
\definecolor{crimson}{RGB}{153, 0, 0}
\definecolor{darkblue}{RGB}{0, 0, 153}

% Hyperlink settings
\hypersetup{
    colorlinks=true,
    linkcolor=crimson,
    urlcolor=crimson,
    citecolor=crimson
}

% Make \textbf default to crimson color
\let\oldtextbf\textbf
\renewcommand{\textbf}[1]{\oldtextbf{\textcolor{crimson}{#1}}}

% Allow line breaks at slashes to prevent overfull hbox
\tolerance=2000
\emergencystretch=2em

% Chapter/section title formatting
\usepackage{titlesec}
\titleformat{\chapter}{\normalfont\huge\bfseries}{}{0pt}{}[\vspace{0.3em}]
\titlespacing*{\chapter}{0pt}{0pt}{14pt}
\titleformat{\section}{\normalfont\Large\bfseries}{\thesection}{1em}{}
\titleformat{\subsection}{\normalfont\large\bfseries}{\thesubsection}{1em}{}

\begin{document}

\pagestyle{plain}

% ─────────────────────────────────────────────
%  Cover Page
% ─────────────────────────────────────────────
\begin{titlepage}
\centering
\vspace*{3cm}

{\fontsize{28}{34}\selectfont\bfseries DSP in VLSI\par}
\vspace{0.6em}
{\fontsize{20}{26}\selectfont\bfseries Homework 2\par}
\vspace{0.4em}
{\fontsize{16}{22}\selectfont\bfseries Digital Filter\par}

\vspace{1.5cm}
\textcolor{crimson}{\rule{0.6\linewidth}{1.2pt}}
\vspace{1.5cm}

{\large\bfseries Student ID}\par
\vspace{0.4em}
{\large M11407439}\par

\vspace{0.8em}

{\large\bfseries Name}\par
\vspace{0.4em}
{\large MING-HONG, LIN}\par

\vspace{2cm}

{\large\bfseries Development Environment}\par
\vspace{0.8em}
\begin{tabular}{rl}
    \textbf{HDL}              & Verilog HDL \\[4pt]
    \textbf{Simulation Tools} & Synopsys VCS and Verdi \\[4pt]
    \textbf{Synthesis Tool}   & Synopsys DC Compiler \\[4pt]
    \textbf{Process}          & TSMC 16nm (ADFP) \\[4pt]
    \textbf{Verification}     & Python \\
\end{tabular}

\vfill

\textcolor{crimson}{\rule{0.6\linewidth}{0.6pt}}\par
\vspace{0.5em}
{\small Due Date: 2026/04/08}

\end{titlepage}

% ─────────────────────────────────────────────
%  Table of Contents
% ─────────────────────────────────────────────
\tableofcontents
\newpage

\chapter{Step 1 — Filter Analysis}

Please use Matlab/Python to draw the \textbf{impulse response} and
\textbf{frequency response} (containing magnitude and phase) of the 25-tap
sinc-based FIR filter. Use dB scale for magnitude and radian for phase versus
normalized frequency. The x-axis must be marked with the correct label.
\textbf{(20\%)}
Please explain whether the filter is high-pass, band-pass or low-pass and why.
\textbf{(5\%)}

\section{Impulse Response}

\begin{figure}[h]
    \centering
    \includegraphics[width=\linewidth]{image0.png}
    \caption{Impulse Response Diagram}
    \label{fig:placeholder}
\end{figure}

\newpage
\section{Frequency Response and Filter Type}

\textbf{This is a low-pass filter.} The filter coefficients are derived from the Sinc function, which corresponds to a rectangular frequency response (ideal low-pass filter) in the frequency domain. From the magnitude response, the passband is located at low frequencies ($\omega<0.141\pi$) with approximately 16 dB gain, and the stopband attenuates high-frequency components by 20 to 40 dB. The 3dB cutoff frequency is at $\omega=0.141\pi$, close to the theoretical value of$\pi/6$. The phase response is linear in the passband, consistent with a symmetric-coefficient FIR filter.

\begin{figure}[h]
    \centering
    \includegraphics[width=1\linewidth]{image1.png}
    \caption{Frequency Response Diagram}
    \label{fig:placeholder}
\end{figure}

\chapter{Step 2 - Input and output waveforms}
Draw the time-domain input and output waveforms after feeding 144 input
samples $x[n] = \sin\!\left(-\tfrac{2\pi n}{128}\right) - \cos\!\left(\tfrac{2\pi n}{4}\right)$,
$0 \le n \le 143$, and observe the filtering effect.
\textbf{(20\%)}

\section{Sinc Function Implementation}
\begin{figure}[h]
    \centering
    \includegraphics[width=\linewidth]{image2.png}
    \caption{Time-Domain Input/Output Waveforms}
    \label{fig:placeholder}
\end{figure}

\chapter{Step 3 — Quantization}

Please use \textbf{direct form} to implement the 25-tap FIR filter. Decide the
word-lengths of each data path after additions and multiplications so that the
RMSE error between the fixed-point output and the floating-point output is less
than $2^{-a}$. If $p = 0$, $a = 12$; otherwise $a = 11$, where
$p = \mathrm{mod}(d_1, 2)$ and $d_1$ is the last digit of your student ID.

To show how you determine the word-length, scan the quantization error versus
word-lengths from 9 bits to 20 bits and mark the final word-length settings in
the block diagram. \textbf{(20\%)}

Draw the fixed-point output of the time-domain signal and compare to the
floating-point output. \textbf{(10\% + 10\%)}
Show the frequency response of the fixed-point direct-form FIR filter and
compare the magnitude error in passband (within 3\,dB bandwidth), expressed in
dB. \textbf{(10\% + 5\%)}

\section{Word-Length Determination}

Based on the analysis of fractional word lengths ranging from \textbf{9 to 20 bits}, the minimum word length required to meet the output error threshold was identified. The experimental results indicate that \textbf{the input requires a 13-bit fractional length}. With the input fixed, \textbf{the coefficients require 15 bits}. Following these assignments, \textbf{the multiplication results require 17 bits}, and finally, \textbf{the accumulation requires 17 bits}.

\begin{figure}[h]
    \centering
    \includegraphics[width=0.98\linewidth]{image4.png}
    \caption{Scan the quantization error versus word-lengths from 9 bits to 20 bits}
    \label{fig:placeholder}
\end{figure}


Furthermore, by incorporating the maximum integer word lengths for each stage, the complete fixed-point configurations are determined as follows:

\begin{table}[h]
    \centering
    \begin{tabular}{lcc}
        \toprule
        \textbf{Stage} & \textbf{Fractional Bits (F)} & \textbf{Total Word Length (W)} \\
        \midrule
        Input          & 13 bits & 16 bits \\
        Coefficient    & 15 bits & 17 bits \\
        Multiplication & 17 bits & 20 bits \\
        Accumulation   & 17 bits & 21 bits \\
        \bottomrule
    \end{tabular}
    \caption{Word-length configuration for each stage.}
    \label{tab:wordlength}
\end{table}


\section{Fixed-Point FIR Analysis}

Based on the experimental results, the overall RMSE is $0.00045324$, which is
below the error threshold $2^{-11} = 0.00048828$. The maximum error in the
passband is $0.00054532$\,dB.

\begin{figure}[h]
    \centering
    \includegraphics[width=1\linewidth]{image3.png}
    \caption{Fixed-Point FIR Analysis}
    \label{fig:placeholder}
\end{figure}


\chapter{Step 4 — Direct Form Implementation \& Synthesis}

Print out the behavior timing diagram for the first 20 samples and the last 20
samples \textbf{(10\%)}. Show the errors of hardware outputs and Matlab
floating-point results of direct form FIR filter output $y[n]$ for
$0 \le n \le 143$ by figures.
\textbf{The latency is not constrained. (10\%)}

Paste the timing report about the critical path. Show the numbers of adders
and multipliers in the critical path \textbf{(5\%)}. List your timing
constraint to achieve the highest operating frequency \textbf{(5\%)}. Print out
the \textbf{post-SYNTHESIS} timing diagram for the first 20 samples and the
last 20 samples \textbf{(10\%)}. and show the errors of post-synthesis
simulation outputs and Matlab floating-point results of direct-form FIR filter
output $y[n]$ for $0 \le n \le 143$ \textbf{(10\%)}.

\section{Behavior Simulation and Error Analysis}

\begin{figure}[h]
    \centering
    \includegraphics[width=1\linewidth]{image6.png}
    \caption{First Timing Diagram (First 20)}
    \label{fig:placeholder}
\end{figure}

\begin{figure}[h]
    \centering
    \includegraphics[width=1\linewidth]{image5.png}
    \caption{Behavior Timing Diagram (Last 20)}
    \label{fig:placeholder}
\end{figure}

\newpage

In the testbench, the filter output is written to \textbf{output.dat}. The
output is then converted from Q4.17 fixed-point to FP32 using Python, and
compared against the original floating-point results for error analysis. The
resulting RMSE is only $0.00045324$, which is below the error threshold.

\begin{figure}[h]
    \centering
    \includegraphics[width=1\linewidth]{image7.png}
    \caption{The errors of hardware outputs and Python floating-point result}
    \label{fig:placeholder}
\end{figure}

\newpage

\section{Critical Path Report and Timing 
Constraint}

The critical path of the direct-form FIR filter consists of \textbf{1 multiplier and 24 adders (Non-Optimization).} The maximum delay from the synthesis timing report is 3.40 ns and \textbf{the maximum operation frequency is 200MHz}.

\begin{figure}[h]
    \centering
    \includegraphics[width=0.45\linewidth]{image12.png}
    \caption{Timing Report}
    \label{fig:placeholder}
\end{figure}

\begin{table}[h]
    \centering
    \begin{tabular}{lc}
        \toprule
        \textbf{Parameter} & \textbf{Value (ns)} \\
        \midrule
        Clock Period                  & 4   \\
        Clock Uncertainty             & 0.5 \\
        Clock Transition              & 0.1 \\
        Clock Latency                 & 0.1 \\
        Input / Output Maximum Delay  & 0.2 \\
        Input Transition              & 0.2 \\
        \bottomrule
    \end{tabular}
    \caption{Timing constraint settings.}
    \label{tab:timing-constraint}
\end{table}

\newpage
\section{Post-Synthesis Simulation and Error Analysis}
Actually, the RMSE of Post-Synthesis is same as Behavior Simulation in Section 3.1 (only $0.00045324$).
\begin{figure}[h]
    \centering
    \includegraphics[width=1\linewidth]{image9.png}
    \includegraphics[width=1\linewidth]{image10.png}
    \caption{Post-Synthesis Timing Diagram (First 20 \& Last 20)}
    \label{fig:placeholder}
\end{figure} 

\begin{figure}[h]
    \centering
    \includegraphics[width=0.75\linewidth]{image11.png}
    \caption{The errors of Post-Synthesis hardware outputs and Python floating-point result}
    \label{fig:placeholder}
\end{figure}

\chapter{Step 5 — 2X Speedup via Parallel Processing}

Show that \textbf{$y[0]$ to $y[143]$ can be generated within
$T_{\mathrm{out}}/2$} in the behavior timing diagram. Draw the block diagram
and indicate your critical path \textbf{(15\%)}. Show your synthesis report of
max delay to verify your explanations \textbf{(5\%)}.

Show the \textbf{post-SYNTHESIS} simulation results for the first 20 samples
and the last 20 samples and indicate why your throughput is doubled
\textbf{(10\%)}. Show the errors of post-synthesis simulation outputs and Matlab
floating-point results of direct-form FIR filter output $y[n]$ for
$0 \le n \le 143$ \textbf{(10\%)}. Show the area (resources) before and after
speedup from the synthesis report \textbf{(10\%)}.

\section{Behavior Simulation}

Experimental results demonstrate that after implementing the parallel processing architecture, the total time required to output the 144th sample $y[n]$ is reduced to \textbf{720 ns (from 85 ns to 805 ns)}. In contrast, the original architecture requires \textbf{1440 ns (from 85 ns to 1525 ns)} to complete the same operation. These data points confirm that the parallel processing architecture effectively halves the processing time, achieving a 2x speedup.

\begin{figure}[h]
    \centering
    \includegraphics[width=1\linewidth]{image14.png}
    \caption{Original Behavior Simulation}
    \label{fig:placeholder}
\end{figure}

\begin{figure}[h]
    \centering
    \includegraphics[width=1\linewidth]{image13.png}
    \caption{2X Speedup Behavior Simulation}
    \label{fig:placeholder}
\end{figure}

\newpage

\section{Critical Path Analysis \& Timing Report}

The design implements a pipeline stage at 5-tap intervals. Consequently, \textbf{the critical path delay consists of 1 multiplier and 4 adders}, significantly reducing the logic depth compared to the non-pipelined baseline.

\begin{figure}[h]
    \centering
    \includegraphics[width=\linewidth]{Filter.png}
    \caption{Block Diagram of 2X Speedup Filter Critical Path}
    \label{fig:placeholder}
\end{figure}

The timing report validates that the critical path is strictly constrained between the successive launch register (PIPE\_STAGE1\_1) and capture register (PIPE\_STAGE1\_2). The total combinational delay is measured at 3.50ns, which accounts for the logic depth of one multiplier and four adders and the timing constraint settings are the same as in Section 4.2.

\begin{figure}[h]
    \centering
    \includegraphics[width=0.5\linewidth]{image15.png}
    \caption{Timing Report of 2X Speedup}
    \label{fig:placeholder}
\end{figure}

\newpage

\section{Post-Synthesis Simulation and Error Analysis}
Actually, the RMSE of Post-Synthesis is alos same as Behavior Simulation in Section 3.1 (only $0.00045324$).

\begin{figure}[h]
    \centering
    \includegraphics[width=0.9\linewidth]{image16.png}
    \includegraphics[width=0.9\linewidth]{image17.png}
    \caption{Post-Synthesis Timing Diagram (First 20 \& Last 20)}
    \label{fig:placeholder}
\end{figure}

\begin{figure}[h]
    \centering
    \includegraphics[width=0.75\linewidth]{image18.png}
    \caption{The errors of Post-Synthesis hardware outputs and Python floating-point result}
    \label{fig:placeholder}
\end{figure}

\newpage

\section{Resources Comparison}

\begin{table}[h]
    \centering
    \begin{tabular}{lcc}
        \toprule
        \textbf{Design} & \textbf{Total Area} & \textbf{Area Overhead} \\
        \midrule
        Original (1$\times$)              & 53{,}504  & ---          \\
        Two-Level Parallel (2$\times$)    & 102{,}624 & 1.92$\times$ \\
        \bottomrule
    \end{tabular}
    \caption{Area comparison before and after 2X speedup.}
    \label{tab:area-comparison}
\end{table}

\end{document}
